\section*{Chapitre \ref{ch:b3:nervous-systems} --- L'organisation des systèmes nerveux}

\begin{solution}{exo:b3:nervous-systems:1}
\omterm{def:b3:nervous-systems:cells}{Astrocytes}~: tamponnent le potassium et le glutamate, nourrissent les
neurones, induisent la \omterm{prop:b3:nervous-systems:barrier-budget}{barrière hémato-encéphalique}, règlent le débit
sanguin local. \omterm{def:b3:nervous-systems:cells}{Oligodendrocytes}~: myélinisent les axones du système
nerveux central. Cellules de Schwann~: myélinisent les axones
périphériques (et guident leur régénération). \omterm{def:b3:nervous-systems:cells}{Microglie}~: les
\omterm{def:b3:innate-immunity:innate}{macrophages} du cerveau, qui élaguent les synapses et nettoient les
débris. (Les cellules épendymaires tapissent les ventricules et
fabriquent le \omterm{prop:b3:nervous-systems:barrier-budget}{liquide céphalorachidien}.)
\end{solution}

\begin{solution}{exo:b3:nervous-systems:2}
L'imprégnation noircissait entièrement quelques neurones et laissait les
autres invisibles, de sorte qu'une cellule unique pouvait être vue en
entier. Cajal en conclut que les neurones sont des cellules séparées qui
se touchent sans fusionner, avec un sens d'écoulement fixe des dendrites
vers l'axone~; Golgi en conclut qu'ils formaient un réseau continu.
Cajal avait raison~; le microscope électronique montra la fente
synaptique en 1954.
\end{solution}

\begin{solution}{exo:b3:nervous-systems:3}
Un coup étire le quadriceps~; ses fuseaux neuromusculaires déchargent~;
les afférents Ia entrent dans la moelle et font synapse directement sur
les motoneurones du quadriceps, qui déchargent et contractent le muscle,
tandis qu'un \omterm{def:b3:nervous-systems:cells}{interneurone} inhibe les motoneurones des ischio-jambiers.
Une synapse, des axones courts à \qty{80}{m/s}~: environ \qty{30}{ms}.
Un coup de pied volontaire passe par le cerveau --- perception,
décision, planification motrice, faisceaux descendants, plusieurs
synapses --- et prend \qtyrange{200}{300}{ms}.
\end{solution}

\begin{solution}{exo:b3:nervous-systems:4}
Sympathique (noradrénaline à la cible)~: cœur plus rapide et plus fort,
pupille dilatée, motricité et sécrétion intestinales réduites, voies
aériennes dilatées. Parasympathique (acétylcholine)~: cœur ralenti,
pupille resserrée, motricité et sécrétion intestinales augmentées, voies
aériennes resserrées.
\end{solution}

\begin{solution}{exo:b3:nervous-systems:5}
$\lambda = \sqrt{R_{m}d/4R_{i}} = \sqrt{10^{4}\times 10^{-4}/400}$ cm
$= \qty{0.05}{cm} = \qty{0.5}{mm}$. Après \qty{2}{mm}~: $e^{-4} \approx
2\,\%$. Doubler $\lambda$ demande quatre fois le diamètre,
\qty{4}{\micro m}.
\end{solution}

\begin{solution}{exo:b3:nervous-systems:6}
Sensoriel~: $12\times 6 = \qty{72}{m/s}$, $1.2/72 = \qty{17}{ms}$.
Moteur~: \qty{90}{m/s}, $1.1/90 = \qty{12}{ms}$. Deux synapses,
\qty{1}{ms}~: environ \qty{30}{ms}. Amyélinique~:
$1.1\sqrt{12} = \qty{3.8}{m/s}$ et
$1.1\sqrt{15} = \qty{4.3}{m/s}$~: $315 + 258 + 1 \approx \qty{570}{ms}$
--- trop lent pour sauver le pied.
\end{solution}

\begin{solution}{exo:b3:nervous-systems:7}
$2.4\times 10^{20}/8.6\times 10^{10} \approx 2.8\times 10^{9}$ ATP par
neurone et par seconde~; à $2\times 10^{9}$ par potentiel, environ
\qty{1.4}{Hz} en moyenne. Le code doit être clairsemé~: peu de
potentiels, l'information portée par les neurones qui déchargent et par
le moment où ils le font, non par des fréquences élevées partout.
\end{solution}

\begin{solution}{exo:b3:nervous-systems:8}
Deux neurones qui s'inhibent l'un l'autre sans se fatiguer forment un
interrupteur~: le premier à décharger étouffe l'autre aussi longtemps
que dure la commande, et le relais ne se passe jamais. L'oscillation
demande un processus lent qui desserre la prise du gagnant ---
adaptation de sa décharge, dépression de sa synapse inhibitrice, ou
rebond du perdant après sa libération --- et dont la constante de temps
fixe la période.
\end{solution}

\begin{solution}{exo:b3:nervous-systems:9}
La L-dopa est transportée à travers la barrière par le transporteur des
acides aminés neutres volumineux~; la dopamine, polaire et chargée, ne
l'est pas. Les patients prennent de la L-dopa (avec un inhibiteur de
l'enzyme périphérique pour qu'elle ne soit pas convertie avant le
passage), et l'enzyme du cerveau lui-même en fait de la dopamine là où
il en faut. Les \omterm{def:b3:adaptive-immunity:antibody}{anticorps}, à \qty{150}{kDa}, ne franchissent ni les
jonctions serrées ni aucun transporteur~; il faut pour les livrer une
navette qui détourne un récepteur pour la transcytose, ou une injection
dans le \omterm{prop:b3:nervous-systems:barrier-budget}{liquide céphalorachidien}.
\end{solution}

\begin{solution}{exo:b3:nervous-systems:10}
Le courant membranaire par unité de longueur est le courant résistif
$V/r_{m}$ plus le courant capacitif $c_{m}\,\partial V/\partial t$~; la
conservation de la charge donne
$\partial I/\partial x = -(V/r_{m} + c_{m}\,\partial V/\partial t)$, et
avec $\partial V/\partial x = -r_{i}I$, $\frac{1}{r_{i}}\,\partial^{2}V/
\partial x^{2} = V/r_{m} + c_{m}\,\partial V/\partial t$~; en multipliant
par $r_{m}$~: $\lambda^{2}\,\partial^{2}V/\partial x^{2} = V + \tau\,\partial
V/\partial t$. Avec $\partial V/\partial t = 0$ l'équation du théorème
revient. $\lambda$ est une longueur et $\tau$ un temps, donc
$\lambda/\tau$ est une vitesse --- l'échelle naturelle de la vitesse à
laquelle une perturbation se propage.
\end{solution}

\begin{solution}{exo:b3:nervous-systems:11}
$v \propto \sqrt{d}$~: quatre fois la vitesse demande seize fois le
diamètre, \qty{8}{mm}. Aire $\pi(4\,\text{mm})^{2} \approx \qty{50}{mm^2}$
contre $\pi\,(10\,\unit{\micro m})^{2} = \qty{314}{\micro m^2}$~: un
rapport de $1.6\times 10^{5}$. Un nerf de fibres amyéliniques rapides est
impossible --- une seule de ces fibres serait aussi épaisse que le nerf
--- de sorte qu'un animal à nombreuses voies rapides a besoin de
\omterm{def:b3:nervous-systems:cells}{myéline}, et c'est pourquoi les vertébrés (et quelques invertébrés,
indépendamment) l'ont fait apparaître.
\end{solution}

\begin{solution}{exo:b3:nervous-systems:12}
$a = 0.4/30^{0.75} = 0.4/12.8 = 0.031$~; pour \qty{70000}{g}~: $0.031\times
\num{70000}^{0.75} = 0.031\times 4300 \approx \qty{134}{g}$. Les
\qty{1350}{g} de l'humain valent dix fois la prévision. La masse du
cerveau suit la masse du corps parce qu'un corps plus grand demande plus
de neurones sensoriels et moteurs, et parce que les neurones eux-mêmes
grossissent avec le cerveau~; ce qui suit la cognition est le nombre de
neurones du cortex, qui dépend de leur densité d'entassement --- les
primates en entassent plus par gramme que les autres mammifères, et les
humains plus que tous.
\end{solution}

\begin{solution}{pb:b3:nervous-systems:1}
\textbf{1.} $\lambda = \sqrt{2\times 10^{4}\,d/400}$~: pour $d =
\qty{1e-4}{cm}$, $\sqrt{5\times 10^{-3}} = \qty{0.071}{cm} =
\qty{0.71}{mm}$~; pour \qty{4}{\micro m}, \qty{1.41}{mm}. $\tau = 2\times
10^{4}\times 10^{-6} = \qty{20}{ms}$.
\textbf{2.} $5\,e^{-0.3/0.71} = \qty{3.3}{mV}$~; $5\,e^{-0.3/1.41} =
\qty{4.0}{mV}$.
\textbf{3.} Les entrées distales arrivent atténuées, si bien que
l'endroit où se pose une synapse fixe son poids~; un potentiel décroît
en $\tau$ environ, si bien que deux entrées ne se somment que si elles
tombent à quelque \qty{20}{ms} l'une de l'autre --- le neurone est un
détecteur de coïncidence à fenêtre de \qty{20}{ms}.
\textbf{4.} Myélinisé \qty{60}{m/s}~; amyélinique $1.1\sqrt{10} =
\qty{3.5}{m/s}$~; pour atteindre \qty{60}{m/s} sans \omterm{def:b3:nervous-systems:cells}{myéline}, $d = (60/1.1)^{2}
\approx \qty{3000}{\micro m}$~: trois millimètres.
\textbf{5.} $\pi(5)^{2} = \qty{79}{\micro m^2}$ contre $\pi(1500)^{2} =
7\times 10^{6}\,\unit{\micro m^2}$~: environ \num{90000} axones
myélinisés dans la place d'un seul.
\textbf{6.} \qty{2}{cm} à \qty{60}{m/s} font \qty{0.3}{ms}~; à
\qty{3.5}{m/s}, \qty{5.7}{ms}~: quelque \qty{5}{ms} de délai
supplémentaire. Pire, la membrane dénudée a une grande capacité et peu
de canaux, si bien que le courant du dernier nœud intact peut ne pas
l'amener au seuil~: le bloc de conduction.
\textbf{7.} Sensoriel $1.2/60 = \qty{20}{ms}$~; deux synapses
\qty{1}{ms}~; moteur $1.1/90 = \qty{12}{ms}$~: environ \qty{33}{ms}
(plus une milliseconde à la jonction neuromusculaire).
\textbf{8.} Vers le cerveau~: $0.6/60 = \qty{10}{ms}$ plus deux synapses,
\qty{11}{ms} après l'entrée du signal dans la moelle à \qty{20}{ms}~: le
signal atteint le cortex à peu près quand le pied bouge~; l'expérience
consciente de la douleur demande quelques centaines de millisecondes de
traitement de plus, de sorte que le pied est retiré avant que la douleur
soit ressentie.
\textbf{9.} $1.1\sqrt{1} = \qty{1.1}{m/s}$~: $1.8/1.1 \approx
\qty{1.6}{s}$. La première douleur, vive, arrive par de fines fibres
myélinisées en quelques dizaines de millisecondes~; la seconde, sourde
et brûlante, par les fibres C amyéliniques une seconde ou plus tard.
\textbf{10.} Sensoriel $3.2/60 = \qty{53}{ms}$, moteur $3.1/90 =
\qty{34}{ms}$, synapses \qty{1}{ms}~: environ \qty{90}{ms}. Les animaux
longs ont des réflexes lents~; ils compensent par des axones plus épais
et en déléguant davantage aux circuits locaux.
\textbf{11.} $0.8/80 = \qty{10}{ms}$ à l'aller comme au retour,
\qty{20}{ms}~; synapse \qty{0.5}{ms}, jonction \qty{1}{ms}, force
\qty{5}{ms}~: \qty{26.5}{ms}, la transduction du fuseau elle-même et la
conduction dans le muscle faisant le reste des \qty{30}{ms}.
\textbf{12.} Les axones amyéliniques conduisent à un ou deux mètres par
seconde, si bien que des boucles qui prennent \qty{30}{ms} chez l'adulte
en prennent des centaines, avec un tempo médiocre~; la myélinisation
multiplie la vitesse par dix à cinquante et rend le tempo précis, ce
qu'exigent la marche, la préhension et la parole.
\textbf{13.} $20/50\,000 = 4\times 10^{-4}$ mol/s $= 2.4\times 10^{20}$
ATP par seconde~; $2.8\times 10^{9}$ par neurone et par seconde.
\textbf{14.} $4\times 10^{8} + 7000\times 2\times 10^{4} = 4\times 10^{8}
+ 1.4\times 10^{8} = 5.4\times 10^{8}$ ATP.
\textbf{15.} Tout aux potentiels~: $2.8\times 10^{9}/5.4\times 10^{8} \approx
\qty{5}{Hz}$~; la moitié~: environ \qty{2.6}{Hz}.
\textbf{16.} À \qty{1}{Hz}, $5.4\times 10^{8}/2.8\times 10^{9} \approx
\qty{20}{\%}$ du budget~: le code est clairsemé --- la plupart des
neurones silencieux la plupart du temps, l'information dans le petit
nombre qui décharge et dans le moment où il le fait.
\textbf{17.} $8.6\times 10^{10}\times 50\times 5.4\times 10^{8} = 2.3
\times 10^{21}$ ATP/s $= 3.9\times 10^{-3}$ mol/s $\approx \qty{190}{W}$
--- le double du corps entier au repos.
\textbf{18.} $5\times 16 = \qty{80}{kJ/h} = \qty{22}{W}$~: tout juste
assez, sans réserve. Quand l'apport s'arrête, les pompes s'arrêtent en
quelques secondes, les gradients se vident, et la conscience se perd en
une dizaine de secondes.
\textbf{19.} $a = 0.4/12.8 = 0.031$~; \qty{70}{kg}~: $0.031\times 4300 =
\qty{134}{g}$~; \qty{5000}{kg}~: $0.031\times 1.06\times 10^{5} \approx
\qty{3300}{g}$.
\textbf{20.} Humain $1350/134 \approx 10$~; éléphant $4800/3300 \approx
1.5$.
\textbf{21.} Dans le \omterm{def:b3:nervous-systems:plan}{cervelet} --- \qty{98}{\%} d'entre eux --- à
coordonner les quarante mille muscles de la trompe. C'est le nombre de
neurones corticaux, non le total, qui suit la cognition~; le \omterm{def:b3:nervous-systems:plan}{cervelet}
suit la tâche motrice de mener un corps énorme.
\textbf{22.} Axones et dendrites s'allongent avec le cerveau, de sorte
que chaque neurone occupe plus de volume et que le nombre de neurones
croît plus lentement que la masse. Chez les primates la taille des
neurones reste presque constante quand le cerveau grossit, si bien que
leur nombre croît presque proportionnellement à la masse, et un cerveau
de primate d'un poids donné contient plusieurs fois les neurones d'un
cerveau de rongeur de ce poids.
\textbf{23.} Un bras peut coordonner ses propres ventouses et
articulations, s'étendre, saisir, faire passer la nourriture le long de
lui-même et éviter d'agripper sa propre peau, par des circuits locaux~;
il ne peut voir, choisir une cible ni apprendre une discrimination.
Test~: couper la chaîne nerveuse du bras du cerveau et toucher le bras
avec de la nourriture --- il la saisit et la fait passer vers l'endroit
où serait la bouche~; un bras amputé en fait autant pendant une heure.
\textbf{24.} Environ $23$ synapses par neurone contre $7000$~: trois
cents fois moins. La carte complète dit quel neurone peut influencer
quel autre, mais non la force ni le signe de chaque connexion, ni quels
neuromodulateurs les changent, ni la dynamique --- de sorte que même
pour le ver le schéma de câblage ne prédit le comportement qu'en partie.
\textbf{25.} $\lambda = \qty{0.71}{mm}$ pour la dendrite de
\qty{1}{\micro m}~; latence du retrait environ \qty{33}{ms}~; fréquence
de décharge moyenne environ \qty{2.6}{Hz} avec la moitié du budget aux
potentiels.
\end{solution}
